% Ссылка в тексте: (приложение~К)
\chapter{КАНВА БИЗНЕС-МОДЕЛИ И ФИНАНСОВАЯ МОДЕЛЬ}

В~приложении представлены канва бизнес-модели по методике Lean~Canvas
и~сценарные расчёты финансовой модели ИИ-методиста
(раздел~4.4 основного текста).

\section*{К.1 — Канва бизнес-модели (Lean Canvas)}

Канва бизнес-модели представлена в~табличной форме.

\begin{table}[htbp]
  \caption{Канва бизнес-модели ИИ-методиста}
  \label{tab:lean-canvas}
  \centering
  \begingroup
  \small
  \setlength{\tabcolsep}{4pt}
  \begin{tabular}{|P{4.1cm}|P{3.9cm}|P{3.9cm}|}
  \hline
  \textbf{Проблема} & \textbf{Решение} & \textbf{Уникальное предложение} \\
  \hline
  Создание 1 часа обучения может занимать 73–490 часов работы методиста;
  высокая стоимость подрядчиков; устаревание контента
  &
  Полуавтоматическая подготовка структуры и~черновых материалов курса
  по~загруженному документу, проактивные подсказки и~ручная проверка
  в~контуре, механизм последующей актуализации
  &
  Сочетание методологической вариативности,
  суверенитета данных (локальный LLM-контур)
  и~глубокой интеграции в~LMS-экосистему
  \\
  \hline
  \textbf{Несправедливое преимущество} & \textbf{Ключевые метрики} & \textbf{Каналы} \\
  \hline
  Переключаемые режимы ADDIE / SAM / BD;
  локальный LLM-контур;
  экосистема LMS Adapstory
  &
  Время до первого методического черновика;
  Покрытие уровней Блума;
  faithfulness;
  себестоимость курса
  &
  Прямые продажи B2B;
  продуктовые демо на~открытом и внутреннем корпусе;
  EdTech-конференции и~контент-маркетинг
  \\
  \hline
  \multicolumn{2}{|P{8.0cm}|}{\centering\textbf{Сегменты потребителей}}
  & \textbf{Структура затрат / доходов} \\
  \hline
  \multicolumn{2}{|P{8.0cm}|}{B2B: HR-отделы, L\&D-команды,
  корпоративные университеты, учебные центры и~EdTech-компании}
  &
  Затраты: ФОТ, продвижение, API-токены, прочие расходы.
  Доходы: B2B-подписка по тарифам 2\,999 / 20\,999 / 44\,999~руб./мес.
  \\
  \hline
  \end{tabular}
  \endgroup
\end{table}

\section*{К.2 — Сценарные финансовые показатели}

Показатели ниже отражают сценарный расчёт. Это не~факт продаж.
Они получены из~помесячного PNL-файла и~используются для проверки
порядка величин: согласуются ли тарифы, расходы и~план привлечения
клиентов с~выбранной архитектурой продукта.
Выручка 2027 и 2028E имеет статус результата выбранных допущений.
Это не traction, не pipeline и не подтверждённая коммерческая воронка.

\begin{table}[htbp]
  \caption{Годовые показатели PNL-модели}
  \label{tab:finance-summary}
  \centering
  \begingroup
  \scriptsize
  \setlength{\tabcolsep}{2pt}
  \renewcommand{\arraystretch}{1.14}
  \begin{tabular}{|c|P{2.25cm}|P{2.35cm}|P{2.55cm}|P{3.55cm}|}
  \hline
  \textbf{Период}
    & \textbf{Выручка}
    & \TableHeadTwo{Расходы}{и налог}
    & \TableHeadThree{Опер.}{денежный}{поток}
    & \textbf{Комментарий} \\
  \hline
  2026 & 76.0 млн руб. & 49.0 млн руб. & 26.9 млн руб. & помесячный расчёт PNL-файла \\
  \hline
  2027 & 1\,764.2 млн руб. & 227.8 млн руб. & 1\,536.4 млн руб. & помесячный расчёт PNL-файла \\
  \hline
  2028E & 4\,044.1 млн руб. & 446.1 млн руб. & 3\,598.0 млн руб. & годовой темп декабря 2027 года \\
  \hline
  Итого & 5\,884.2 млн руб. & 722.9 млн руб. & 5\,161.4 млн руб. & сценарный горизонт 2026–2028E \\
  \hline
  \end{tabular}
  \endgroup
\end{table}

\section*{К.3 — Допущения модели}

Допущения задают границы интерпретации финансового сценария.
При изменении тарифов, конверсии, оттока или ФОТ итоговые значения
требуют пересчёта.

\begin{table}[htbp]
  \caption{Ключевые допущения PNL-модели}
  \label{tab:finance-assumptions}
  \centering
  \begingroup
  \small
  \setlength{\tabcolsep}{4pt}
  \begin{tabular}{|P{4.0cm}|P{8.1cm}|}
  \hline
  \textbf{Блок модели} & \textbf{Допущение} \\
  \hline
  Тарифы & 2\,999 / 20\,999 / 44\,999~руб. в месяц в зависимости от пакета продукта \\
  \hline
  Продажи & три сценария привлечения: органическое продвижение, поисковые запросы, маркетинговая стратегия и SEO \\
  \hline
  Отток & 10\,\% активной базы в расчёте пользовательского потока \\
  \hline
  ФОТ & стартовый ФОТ 1.06~млн~руб. в месяц с последующим ростом по сценарию \\
  \hline
  Маркетинг & базовый бюджет 300~тыс.~руб. с ростом маркетинговой статьи по сценарию \\
  \hline
  Налоги & УСН 6\,\% от выручки; взносы от ФОТ учтены отдельной строкой \\
  \hline
  \end{tabular}
  \endgroup
\end{table}

\begin{itemize}
  \item исходный PNL-файл содержит детальный помесячный расчёт на 2026–2027 годы;
  \item 2028E в таблице рассчитан по годовому темпу
    декабря 2027 года;
  \item модель носит сценарный характер и~не~заменяет данные будущих продаж;
  \item расчёт прежде всего подтверждает экономическую устойчивость
    выбранной архитектуры;
  \item коммерческие выводы не используются для доказательства качества MVP;
    качество подтверждается отдельными техническими замерами на внутреннем
    и открытом корпусе.
\end{itemize}

\section*{К.4 — Чувствительность финансовой модели}

Для защиты финансового блока модель дополнена чувствительностью по трём
сценариям. Она показывает, какие коммерческие допущения сильнее всего
влияют на выручку, LTV/CAC и срок окупаемости.

\begin{table}[htbp]
  \caption{Сценарная чувствительность PNL-модели}
  \label{tab:appendix-finance-sensitivity}
  \centering
  \begingroup
  \scriptsize
  \setlength{\tabcolsep}{1.6pt}
  \renewcommand{\arraystretch}{1.12}
  \begin{tabular}{|P{1.8cm}|P{3.0cm}|P{2.0cm}|P{2.0cm}|P{2.0cm}|P{3.0cm}|}
  \hline
  \textbf{Сценарий}
    & \textbf{Драйверы}
    & \textbf{2027}
    & \textbf{2028E}
    & \textbf{Unit-экономика}
    & \textbf{Условие применимости} \\
  \hline
  Conservative
    & 40\,\% базового плана продаж; CAC выше плана; отток до 15\,\%
    & $\approx$706 млн руб.
    & $\approx$1\,618 млн руб.
    & LTV/CAC 3--4; окупаемость 6--8 мес.
    & найм и маркетинг растут только после подтверждённых оплат \\
  \hline
  Base
    & текущий PNL-план; CAC около 30 тыс. руб.; отток 10\,\%
    & 1\,764.2 млн руб.
    & 4\,044.1 млн руб.
    & LTV/CAC 10; окупаемость 2 мес.
    & повторяемый B2B-канал продаж подтверждает план привлечения \\
  \hline
  Aggressive
    & 130\,\% базового плана продаж; CAC ниже плана; отток до 7\,\%
    & $\approx$2\,293 млн руб.
    & $\approx$5\,257 млн руб.
    & LTV/CAC 13--15; окупаемость 1--2 мес.
    & доказана масштабируемость продаж без кратного роста поддержки \\
  \hline
  \end{tabular}
  \endgroup
\end{table}

\section*{К.5 — Драйверы и условия сходимости}

Ключевые драйверы модели:
\begin{itemize}
  \item платящие B2B-клиенты по месяцам и распределение по тарифам
    2\,999 / 20\,999 / 44\,999~руб.;
  \item CAC, конверсия из демонстрации в оплату и длительность цикла сделки;
  \item отток, продления и фактический срок жизни клиента;
  \item ФОТ, маркетинг, поддержка внедрений и темп найма;
  \item себестоимость генерации курса, API-токены и GPU-стоимость.
\end{itemize}

Что должно быть правдой, чтобы базовая модель сошлась:
\begin{itemize}
  \item CAC удерживается около 30 тыс. руб. без деградации качества лидов;
  \item тарифный микс включает корпоративные пакеты, не только минимальный тариф;
  \item месячный отток удерживается около 10\,\% или ниже;
  \item ФОТ и маркетинговые расходы растут после появления подтверждённой выручки;
  \item инфраструктурная себестоимость одного курса остаётся около 75~руб.;
  \item юридические и интеграционные требования B2B-клиентов не добавляют
    непропорциональную стоимость внедрения.
\end{itemize}
